\subsubsection{Component diagnosis with paired cue flips}
\label{app:exp3-components}

The preceding RL pilot motivates a fixed, exploratory decomposition, frozen
before these inference outputs. We sample 96 fresh latent episodes, 24 in each
of the two domain $\times$ two label cells. Each episode is presented with both
cue assignments under four interfaces, giving 768 prompts per checkpoint:
(i) identify the named reference without forecasting; (ii) forecast with exact
unit-shock response coefficients at horizons 1 and 3 supplied for both
references; (iii) forecast with the selected reference and its coefficients
supplied; and (iv) the original forecasting task. The coefficients are latent
simulator summaries, not the requested ten forecasts; this assistance tests
whether a supplied relationship is usable, not whether exact coefficients are
recoverable from fourteen noisy observations. We use $k=0$ target observations
so numerical evidence is identical across the original-task cue flips. This
diagnostic does not test integration with target observations at $k>0$.

The fixed roster comprises Qwen3-8B base, its matched and fixed-prior RL
checkpoints at seeds 42--44, and an untrained Qwen3-32B feasibility comparison.
All 6,144 greedy generations parse. Selection-only accuracy is $100\%$ for
every checkpoint in every domain/label cell. Table~\ref{tab:exp3-components}
reports the frozen scores, including every assistance condition. At 8B,
supplying patterns does not reduce average response error in any domain/label
cell for the base or either trained arm. This locates a failure beyond merely
identifying the reference named by the cue; it does not uniquely separate
numerical application from sensitivity to the assisted prompt format.

\begin{table}[htbp]
\centering\small
% Generated by paper/generate_exp3_diagnostic_tables.py; do not edit.
\begin{tabular}{lllrrrr}
\toprule
Domain & Labels & Model & Name (\%) & Original & Both & Selected \\
\midrule
City & semantic & 8B base & 100 & 2.105 & 2.393 & 2.413 \\
City & semantic & 8B matched RL & 100 & 1.053 & 2.213 & 1.440 \\
City & semantic & 8B prior RL & 100 & 0.990 & 3.657 & 1.224 \\
City & semantic & 32B base & 100 & 4.110 & 1.710 & 1.423 \\
City & arbitrary & 8B base & 100 & 2.117 & 2.710 & 2.961 \\
City & arbitrary & 8B matched RL & 100 & 1.130 & 3.793 & 1.567 \\
City & arbitrary & 8B prior RL & 100 & 1.084 & 4.156 & 1.342 \\
City & arbitrary & 32B base & 100 & 2.671 & 1.993 & 1.604 \\
Harbor & semantic & 8B base & 100 & 4.967 & 7.841 & 7.872 \\
Harbor & semantic & 8B matched RL & 100 & 3.318 & 9.657 & 8.633 \\
Harbor & semantic & 8B prior RL & 100 & 2.903 & 7.938 & 7.197 \\
Harbor & semantic & 32B base & 100 & 7.543 & 5.288 & 4.730 \\
Harbor & arbitrary & 8B base & 100 & 5.353 & 12.155 & 8.826 \\
Harbor & arbitrary & 8B matched RL & 100 & 3.418 & 13.837 & 9.303 \\
Harbor & arbitrary & 8B prior RL & 100 & 3.015 & 10.971 & 8.224 \\
Harbor & arbitrary & 32B base & 100 & 7.012 & 5.242 & 5.200 \\
\bottomrule
\end{tabular}

\caption{\textbf{Reference naming succeeds while forecasting remains difficult.}
Name is selection-only accuracy; Original, Both, and Selected are response MAE
for the three forecast interfaces. Trained rows average all three seeds, not
selected endpoints. Each checkpoint contributes 24 paired episodes per cell.
Domain scales differ, so compare interfaces within a domain.}
\label{tab:exp3-components}
\end{table}

Forecast structure accuracy assigns each predicted eight-coordinate response
vector to the closer of the episode's direct and mediated simulator templates
in MAE; ties count as incorrect. The paired cue-change calibration is
\[
\alpha=\frac{1}{N}\sum_{i=1}^N
\frac{(\widehat r_{i,M}-\widehat r_{i,D})^\top(r_{i,M}-r_{i,D})}
     {\lVert r_{i,M}-r_{i,D}\rVert_2^2},
\]
where $r$ subtracts the zero-shock forecast separately within each horizon,
and $M,D$ index the mediated and direct cue assignments. Zero denotes
invariance to the cue; one denotes the simulator's change along this direction,
but need not imply a correct full response vector. This differs from the
preceding persistence-ratio index, whose selecting reference is approximately
$.65$. Matched RL's mean $\alpha$ is $.000$, $.017$, $.136$, and $.008$ for
City semantic, City arbitrary, Harbor semantic, and Harbor arbitrary,
respectively. The corresponding fixed-prior means are $-.008$, $.027$,
$.085$, and $.035$; these three-seed descriptions are not a new confirmatory
comparison between training arms.

\paragraph{Stronger-model feasibility and output semantics.}
The frozen scorer clips forecasts to the valid domain range before computing
responses. Inspection revealed that Qwen3-32B often returns deviations from rest
instead of absolute levels; clipping those values distorts the response vector,
especially in Harbor, whose lower bound is 40. We therefore report a separately
labelled \emph{post hoc} raw-output audit (Table~\ref{tab:exp3-baseline-audit}),
without replacing the frozen scores or counting offset-corrected answers as
correct forecasts. With the selected pattern supplied in City/semantic, raw
response MAE is $.384$ versus $4.794$ on the original interface. Complete
relative vectors are accurate to $.01$ in $52.1\%$ of cases, but $95.8\%$ omit
the resting level. This shows partial conditional numerical ability, not a
solution of the original task. The 8B assistance failures remain predominantly
response-pattern errors after a common-offset correction.

\begin{table}[htbp]
\centering\small
% Generated by paper/generate_exp3_diagnostic_tables.py; do not edit.
\begin{tabular}{llrrrrr}
\toprule
Domain & Labels & Original & Both & Selected & Exact (\%) & Omit rest (\%) \\
\midrule
City & semantic & 4.794 & 0.883 & 0.384 & 52.1 & 95.8 \\
City & arbitrary & 2.721 & 1.225 & 0.525 & 41.7 & 91.7 \\
Harbor & semantic & 10.754 & 2.910 & 3.187 & 18.8 & 85.4 \\
Harbor & arbitrary & 7.652 & 1.500 & 2.494 & 35.4 & 85.4 \\
\bottomrule
\end{tabular}

\caption{\textbf{Post hoc Qwen3-32B audit of incorrectly offset answers.}
The first three numeric columns are unclipped response MAE. Exact means all
eight relative response coordinates are within $.01$ of truth; Omit rest means
both zero-shock predictions are near zero. These last two columns concern the
selected-pattern interface only. They do not measure correct absolute forecasts.}
\label{tab:exp3-baseline-audit}
\end{table}

\subsubsection{Supervised learnability and its transfer boundary}
\label{app:exp3-learnability}

We train two fixed forecast-SFT recipes from the same Qwen3-8B base on the
original 4,800 balanced Coin City/semantic prompts, with simulator-correct
absolute forecasts as targets. Both use seed 42, one epoch, effective batch 16,
300 optimizer updates, completion-only token cross entropy, LoRA rank 32 and
alpha 64 on all seven projection modules, AdamW $(.9,.95)$, no weight decay,
constant learning rate, and gradient norm cap 1.0. Rates are $10^{-6}$ (the RL
rate) and $2\times10^{-5}$ (an optimization-feasibility condition). Neither
Harbor nor arbitrary labels appear in training. Final endpoints alone are
evaluated, on the unchanged 768-prompt component diagnostic.

This matches distinct prompts and prompt presentations, not computation or
optimizer updates: RL samples eight completions per prompt and uses about
2,400 updates. The higher-rate SFT recipe also changes learning rate. Both
supervised outcomes are reported in Table~\ref{tab:exp3-learnability}; the
seed-42 RL checkpoint is the primary diagnostic comparison. The other two
matched RL seeds yield City/semantic structure accuracies $52.1\%$ and
$50.0\%$, with response MAE $1.029$ and $1.053$.

\begin{table}[htbp]
\centering\small
\setlength{\tabcolsep}{4pt}
% Generated by paper/generate_exp3_diagnostic_tables.py; do not edit.
\begin{tabular}{lllrrrr}
\toprule
Domain & Labels & Recipe & Structure (\%) & Resp. MAE & $\alpha$ & 95\% interval \\
\midrule
City & semantic & Base & 45.8 & 2.105 & -0.211 & [-0.467, 0.014] \\
City & semantic & RL & 50.0 & 1.077 & 0.005 & [-0.074, 0.078] \\
City & semantic & SFT $10^{-6}$ & 58.3 & 1.253 & 0.166 & [-0.235, 0.549] \\
City & semantic & SFT $2\!\times\!10^{-5}$ & 97.9 & 0.658 & 1.413 & [1.205, 1.620] \\
City & arbitrary & Base & 54.2 & 2.117 & 0.344 & [0.162, 0.529] \\
City & arbitrary & RL & 50.0 & 1.222 & 0.012 & [-0.068, 0.102] \\
City & arbitrary & SFT $10^{-6}$ & 43.8 & 1.727 & 0.001 & [-0.424, 0.489] \\
City & arbitrary & SFT $2\!\times\!10^{-5}$ & 39.6 & 1.274 & -0.211 & [-0.516, 0.107] \\
Harbor & semantic & Base & 43.8 & 4.967 & 0.170 & [-0.157, 0.499] \\
Harbor & semantic & RL & 50.0 & 3.445 & 0.180 & [0.050, 0.325] \\
Harbor & semantic & SFT $10^{-6}$ & 54.2 & 5.317 & 0.058 & [-0.447, 0.538] \\
Harbor & semantic & SFT $2\!\times\!10^{-5}$ & 45.8 & 12.192 & -0.932 & [-1.528, -0.375] \\
Harbor & arbitrary & Base & 50.0 & 5.353 & -0.094 & [-0.301, 0.109] \\
Harbor & arbitrary & RL & 52.1 & 3.619 & 0.026 & [-0.115, 0.190] \\
Harbor & arbitrary & SFT $10^{-6}$ & 47.9 & 5.965 & 0.120 & [-0.132, 0.381] \\
Harbor & arbitrary & SFT $2\!\times\!10^{-5}$ & 50.0 & 10.655 & 0.004 & [-0.282, 0.255] \\
\bottomrule
\end{tabular}

\caption{\textbf{A supervised recipe learns in-domain structure selection but
fails to transfer.} All trained recipes here use seed 42. Structure accuracy
is derived from forecasts, not reference naming. All outputs parse. Intervals
resample the 24 paired episodes within each cell (2,000 draws); they do not
capture training-seed uncertainty. Both fixed SFT learning rates are retained.}
\label{tab:exp3-learnability}
\end{table}

At the higher rate, City/semantic structure accuracy is $47/48$ ($97.9\%$),
response MAE is $.658$, and forecast-level MAE is $.527$. Calibration
$\alpha=1.413$ indicates overshoot despite accurate choice of response form.
The lower-rate recipe reaches only $58.3\%$ structure accuracy. Higher-rate
SFT does not recover structure selection under arbitrary labels or in Harbor;
its Harbor response errors increase substantially. Original-task SFT response
scores are unchanged when clipping is removed, and an additive-offset audit
does not remove the transfer failure. Reference naming remains perfect.
Thus the successful recipe establishes in-domain learnability, while a
semantic cue-to-pattern association remains a possible explanation. It does
not establish transferable induction from reference trajectories or an isolated
effect of SFT versus RL.

\paragraph{Adaptive decision disclosure.}
The initially specified numerical launch gate was not met. Before any SFT
training, a documented outcome-informed amendment used the stronger model's
conditional numerical ability and absolute-level errors to justify running the
two already prepared, unchanged recipes. This is an exploratory learnability
pilot, not a passed preregistered gate or a confirmatory method comparison.
There was no SFT endpoint selection. Separate reference-identification
supervision was unnecessary given the perfect selection-only results.

\subsubsection{Mechanism-family evidence use under coherent gain changes}
\label{app:exp3-evidence-use}

The mechanism-family accuracy comparison uses a fixed population-prior training
control. It is distinct from the parent Coin City shuffled-pair control and
does not preserve the matched arm's exact target-vector distribution. To test
whether the accuracy gain depends on displayed episode-specific relationships,
we freeze a fresh gain-intervention set before inference: 16 pairs in each of
12 mechanisms, comprising 128 pairs from eight seen mechanisms and 64 from
four held-out compositions or parameter-extrapolation mechanisms. Both
mechanism-disclosure conditions use the same episodes, with $k=9$ target
observations.

Within each pair the target gain changes from the 20th to the 80th percentile
of that world's native uniform support. Inputs, exact noise draws, calibration
trajectories, baseline, and mechanism are fixed; affected target observations,
terminal states, and future outcomes are regenerated coherently. The disclosed
description provides the same gain range on both sides, never the actual gain.
This tests a gain change, not a change of mechanism identity. Native supports
are retained, including the narrow, noisy extrapolation world.

The primary metric is response-change MAE: subtract the zero-shock forecast
within each horizon, take the high-minus-low change of the resulting eight
responses, and compare with the simulator's change. This removes purely
additive level tracking. Greedy inference evaluates Qwen3-4B-Instruct-2507
base and all retained final matched/fixed-prior adapters, seeds 45 and 46;
weights for seeds 42--44 are unavailable. All pairs parse. The two-seed mean
contrasts and bootstrap intervals in Table~\ref{tab:exp3-gain-evidence} resample
paired training seeds and world-stratified episodes (2,000 draws). With only
two training seeds these are exploratory intervals, not seed-generalized
confirmation. Individual seed results are retained in the released analysis.

\begin{table}[htbp]
\centering\footnotesize
\setlength{\tabcolsep}{3pt}
% Generated by paper/generate_exp3_diagnostic_tables.py; do not edit.
\begin{tabular}{llrrrrrrrr}
\toprule
Disclosure & Mechanisms & Base & Matched & Shuffled & No change & Slope M & Slope S & Shuffled$-$matched & 95\% interval \\
\midrule
Disclosed & Seen & 3.405 & 1.250 & 1.512 & 1.591 & 0.323 & 0.084 & +0.262 & [+0.212, +0.319] \\
Disclosed & Held out & 2.629 & 0.898 & 0.874 & 0.826 & 0.122 & 0.049 & -0.024 & [-0.079, +0.027] \\
Undisclosed & Seen & 3.691 & 1.115 & 1.488 & 1.591 & 0.461 & 0.107 & +0.373 & [+0.318, +0.426] \\
Undisclosed & Held out & 3.664 & 0.878 & 0.862 & 0.826 & 0.238 & 0.060 & -0.016 & [-0.061, +0.027] \\
\bottomrule
\end{tabular}

\caption{\textbf{Better gain tracking on seen mechanisms does not establish
transferred evidence use.} Entries are response-change MAE (lower is better);
trained entries average seeds 45 and 46. No change is an invariant predictor.
Positive prior-minus-matched contrasts favor matched training. Both held-out
intervals include zero, and matched held-out errors exceed the no-change baseline.}
\label{tab:exp3-gain-evidence}
\end{table}

Matched training improves on the prior for seen mechanisms by $.195$
(disclosed) and $.386$ (undisclosed), but held-out contrasts are $-.019$ and
$.007$. Held-out matched tracking slopes are $.102$ and $.233$ (ideal 1),
and its errors exceed the $.826$ no-change baseline in both conditions.
The previously reported ordinary forecast-accuracy advantage therefore remains
a narrower result than transferred evidence use. These interventions neither
show a complete absence of sensitivity nor establish robust tracking of
held-out episode-specific gains.

\paragraph{Control status and reproducibility.}
No completed distribution-matched shuffled-target training comparison is
reported for the mechanism-family benchmark. The separately frozen design
preserves identical prompts and the exact joint target-vector multiset within
compatible world/evidence-depth groups, changing only episode assignments,
with five seeds in each disclosure condition. Only completed diagnostics enter the tables and claims here. Their artifacts,
manifests, raw
outputs, and the outcome-informed SFT amendment reside under
\path{exp3_training_transfer/overnight_diagnostics_20260921/}, in
\path{structure/}, \path{learnability/}, and \path{evidence_use/}.
\path{paper/generate_exp3_diagnostic_tables.py} renders these four tables from
completed summaries and records source hashes in
\path{paper/tables/exp3_diagnostics_sources.json}.
